\section{Roster Results}
\label{sec:results}

Table~\ref{tab:pooled} reports the roster on the frozen test set, fitted on
training plus development and scored once.

\begin{table}[t]
\caption{Pooled test scores, best imbalance arm per model. Intent and priority
are \macrof{}, sentiment \negf{}. Dashes mark cells not yet run.}
\label{tab:pooled}
\centering
\small
\begin{tabular}{lrrr}
\toprule
Model & Intent & Sentiment & Priority \\
\midrule
LaBSE                  & \textbf{0.8835} & \textbf{0.7138} & 0.8900 \\
XLM-R base             & 0.8801 & 0.7007 & 0.8872 \\
mmBERT                 & 0.8680 & 0.7000 & 0.8887 \\
MuRIL                  & 0.8467 & 0.6790 & 0.8717 \\
Gemma 3 1B             & 0.8635 & 0.7126 & 0.8898 \\
Gemma 3 270M           & 0.8414 & --     & --     \\
TF-IDF + linear SVM    & 0.8308 & 0.6653 & 0.8734 \\
TF-IDF + logistic      & 0.8189 & 0.6383 & 0.8706 \\
TF-IDF + complement NB & 0.6792 & 0.4968 & 0.8422 \\
LaBSE frozen probe     & 0.7822 & --     & 0.8015 \\
\midrule
\emph{Label ceiling}   & n/a & \emph{0.7931} & \emph{0.7722} \\
\bottomrule
\end{tabular}
\end{table}

Three readings are supported and a fourth is not.

\textbf{The encoder advantage is real on every task and largest on sentiment.}
LaBSE beats the linear SVM by $+0.0485$ \negf{} (95\% CI $[0.026, 0.072]$,
McNemar $p=1.4\times10^{-7}$) and by $+0.0526$ \macrof{} on intent (CI
$[0.048, 0.058]$, $p=3.2\times10^{-87}$). Section~\ref{sec:script} shows this
pooled figure hides where the gain lives.

\textbf{Sentiment is a tie at the top and we name no winner.} LaBSE, Gemma 3 1B,
and two multitask variants span 0.0096, while a single seed repeat of LaBSE moved
it by 0.0158 at matched budget, larger than the spread it would decide. The
defensible statement is that every encoder and decoder lands at $0.70\pm0.01$
against a classical floor of 0.665, and that the best reaches 90.0\% of what the
labels themselves support.

\textbf{A 270M decoder buys nothing over character $n$-grams on intent},
scoring 0.8414 against 0.8308 at far higher training cost, with both five points
below LaBSE. At this corpus size, parameter count is not what separates systems.

\textbf{What the table does not support} is a decoder scaling claim as
originally run, since the 1B model fell back to batch 8 after an out-of-memory
failure. Re-running the 270M model at batch 8 leaves parameter count as the only
difference and gives $+0.0221$.
